\documentclass[sn-mathphys,Numbered]{sn-jnl}

%% Standard packages (you can add more later if needed)
\usepackage{graphicx}
\usepackage{amsmath,amssymb,amsfonts}
\usepackage{booktabs}
\usepackage{xcolor}

\begin{document}

%% -----------------------------------------------------------------
%% Title and authors
%% -----------------------------------------------------------------

\title[Smart Agriculture Prediction System]{Smart Agriculture Prediction System\\[4pt]
\small Supervisor: PhD. Emiliano Mankolli}

\author*[1]{\fnm{Brajan} \sur{Gjorga}}
\author[1]{\fnm{Ersa} \sur{M\"ezuraj}}
\author[1]{\fnm{Ersona} \sur{Duro}}
\author[1]{\fnm{Kanita} \sur{Kadiu}}

\affil*[1]{\orgdiv{Department of Engineering and Architecture}, \orgname{University of New York Tirana}, \city{Tirana}, \country{Albania}}

%% -----------------------------------------------------------------
%% Abstract and keywords (you can overwrite later)
%% -----------------------------------------------------------------

\abstract{This report presents the design and evaluation of the \textit{Smart Agriculture Prediction System}, a machine-learning based approach for controlling irrigation and ventilation in a greenhouse environment.}

\keywords{smart agriculture, irrigation, prediction system, machine learning}

\maketitle

\clearpage
\tableofcontents
\clearpage

%% -----------------------------------------------------------------
%% Main sections (placeholders for your text)
%% -----------------------------------------------------------------

\section{Introduction}\label{sec:intro}


\section{Introduction}\label{sec:intro}

 on the real work you did in \texttt{Irrigation\_TimeSeries.ipynb}:
% - Smart agriculture context and problem definition
% - Why we need predictions for water pump and fan (ON/OFF)
% - High-level description of the time-aware model and the demo web app


\section{Methodology}\label{sec:methodology}

% TODO: Describe the methodology here, based on Irrigation_TimeSeries.ipynb and the Flask app.
% Suggested points:
% - Dataset (sensors: temperature, humidity, water level, N, P, K; time stamps; actuator states)
% - Cleaning and target construction (irrigation and fan ON/OFF)
% - Feature engineering: month, hour, day_of_week, is_weekend, nutrient_index
% - Time-based train/test split (no shuffling)
% - Model: multi-output RandomForest in a sklearn Pipeline (preprocessor + classifier)
% - Thresholds (0.85, 0.9, 0.95) and the choice of 0.9 for the app
% - 10\% noise robustness experiment


\section{Results}\label{sec:results}

% TODO: Summarize the main experimental results.
% Suggested points:
% - Metrics for irrigation and fan prediction on the time-based test set
% - Confusion matrices for thresholds 0.85, 0.9, 0.95
% - Effect of adding 10\% noise to training features
% - Feature importance discussion (which sensors and time-features matter most)
% - Short description of how the demo app uses the pipeline (``what to do now'' + 24h schedule)


\section{Conclusion}\label{sec:conclusion}

% TODO: Write the conclusion here.
% Suggested points:
% - What the Smart Agriculture Prediction System achieves in practice
% - Benefits and limitations of the current model and app
% - Ideas for future improvements (better sensors, more advanced time-series models, richer schedules)


\end{document}

